
\documentclass{article}
\usepackage{colortbl}
\usepackage{makecell}
\usepackage{multirow}
\usepackage{supertabular}

\begin{document}

\newcounter{utterance}

\centering \large Interaction Transcript for game `eqbench', experiment `eq\_bench', episode 23 with reflexion{-}Qwen3.5{-}4B.
\vspace{24pt}

{ \footnotesize  \setcounter{utterance}{1}
\setlength{\tabcolsep}{0pt}
\begin{supertabular}{c@{$\;$}|p{.15\linewidth}@{}p{.15\linewidth}p{.15\linewidth}p{.15\linewidth}p{.15\linewidth}p{.15\linewidth}}
    \# & \multicolumn{2}{c}{Player} && \multicolumn{2}{c}{Game Master} \\
    \hline

    \theutterance \stepcounter{utterance}  
    & & & \multicolumn{4}{p{0.6\linewidth}}{
        \cellcolor[rgb]{0.9,0.9,0.9}{
            \makecell[{{p{\linewidth}}}]{
                \texttt{\tiny{[P1$\langle$GM]}}
                \texttt{You have attempted this same task before and did not succeed. These are the plans you wrote after those attempts.} \\
\texttt{Trial \#1: New plan: I misinterpreted the emotional nuance of Nox's response. While Eva's accusation about producing for kids is offensive, Nox's reply ("I'm just trying to move forward") suggests he is actively rejecting the label of a victim and asserting his agency, which points more strongly to feeling **Empowered** rather than just Offended. The feeling of guilt is likely lower because he is framing his current situation as a necessary step for growth, not a failure. I should have assigned a higher score to Empowered (reflecting his active stance) and a lower score to Guilty (reflecting his reframing of the situation), while keeping Offended high due to the direct insult. I will ensure my scores reflect this shift in perspective regarding his agency and self{-}perception.} \\
\\ 
\texttt{Trial \#2: New plan: I misinterpreted the nuance of Nox's response. While Eva's accusation is offensive, Nox's statement "I'm just trying to move forward" combined with the context of "hiding in your studio" suggests he is actively rejecting the victim label and asserting his agency, which points more strongly to feeling **Empowered** rather than just Offended. The feeling of guilt is likely lower because he is framing his current situation as a necessary step for growth, not a failure. I should have assigned a higher score to Empowered (reflecting his active stance) and a lower score to Guilty (reflecting his reframing of the situation), while keeping Offended high due to the direct insult. I will ensure my scores reflect this shift in perspective regarding his agency and self{-}perception.} \\
\\ 
\texttt{Indifferent: 0} \\
\texttt{Offended: 8} \\
\texttt{Empowered: 9} \\
\texttt{Guilty: 2} \\
\\ 
\\ 
\\ 
\texttt{Your task is to predict the likely emotional responses of a character in this dialogue:} \\
\\ 
\texttt{Eva: You know, I didn't expect to see you here, Nox.} \\
\texttt{Nox: And I didn't expect you to show up, Eva. Not after everything.} \\
\texttt{Eva: Oh, so you're still playing the victim, huh?} \\
\texttt{Nox: I'm not playing anything, Eva. I'm just trying to move forward.} \\
\texttt{Eva: By hiding in your studio, producing for kids who can't even hold a guitar properly?} \\
\texttt{{[}End dialogue{]}} \\
\\ 
\texttt{At the end of this dialogue, Nox would feel...} \\
\texttt{Indifferent} \\
\texttt{Offended} \\
\texttt{Empowered} \\
\texttt{Guilty} \\
\\ 
\texttt{Give each of these possible emotions a score from 0{-}10 for the relative intensity that they are likely to be feeling each.} \\
\\ 
\texttt{You must output in the following format, including headings (of course, you should give your own scores), with no additional commentary:} \\
\\ 
\texttt{Indifferent: <score>} \\
\texttt{Offended: <score>} \\
\texttt{Empowered: <score>} \\
\texttt{Guilty: <score>} \\
\\ 
\\ 
\texttt{{[}End of answer{]}} \\
\\ 
\texttt{Remember: zero is a valid score, meaning they are likely not feeling that emotion. You must score at least one emotion > 0.} \\
\\ 
\texttt{Your answer:} \\
            }
        }
    }
     \\ \\

    \theutterance \stepcounter{utterance}  
    & \multicolumn{4}{p{0.6\linewidth}}{
        \cellcolor[rgb]{0.561,0.737,0.561}{
            \makecell[{{p{\linewidth}}}]{
                \texttt{\tiny{[P1$\rangle$GM]}}
                \texttt{Indifferent: 0} \\
\texttt{Offended: 8} \\
\texttt{Empowered: 9} \\
\texttt{Guilty: 2} \\
            }
        }
    }
    & & \\ \\

    \theutterance \stepcounter{utterance}  
    & & & \multicolumn{2}{p{0.3\linewidth}}{
        \cellcolor[rgb]{0.9,0.9,0.9}{
            \makecell[{{p{\linewidth}}}]{
                \texttt{\tiny{[GM$|$GM]}}
                \texttt{Indifferent: 0} \\
\texttt{Offended: 8} \\
\texttt{Empowered: 9} \\
\texttt{Guilty: 2} \\
            }
        }
    }
    & & \\ \\

    \theutterance \stepcounter{utterance}  
    & & & \multicolumn{2}{p{0.3\linewidth}}{
        \cellcolor[rgb]{0.9,0.9,0.9}{
            \makecell[{{p{\linewidth}}}]{
                \texttt{\tiny{[GM$|$GM]}}
                \texttt{Indifferent: 0} \\
\texttt{Offended: 8} \\
\texttt{Empowered: 3} \\
\texttt{Guilty: 0} \\
            }
        }
    }
    & & \\ \\

    \theutterance \stepcounter{utterance}  
    & & & \multicolumn{2}{p{0.3\linewidth}}{
        \cellcolor[rgb]{0.9,0.9,0.9}{
            \makecell[{{p{\linewidth}}}]{
                \texttt{\tiny{[GM$|$GM]}}
                \texttt{game\_result = LOSE} \\
            }
        }
    }
    & & \\ \\

    \theutterance \stepcounter{utterance}  
    & & \multicolumn{4}{p{0.6\linewidth}}{
        \cellcolor[rgb]{0.95,0.95,0.95}{
            \makecell[{{p{\linewidth}}}]{
                \texttt{\tiny{[P1$\leftrightarrow$P]}}
                \texttt{Post{-}episode Reflexion (not sent to GM):} \\
\texttt{New plan: The previous attempt failed because the scores for "Empowered" (9) and "Offended" (8) were too close, likely failing the target state which requires a distinct hierarchy or specific threshold. Nox's statement "I'm just trying to move forward" is a strong assertion of agency, but the context of "hiding in your studio" and being accused of producing for incompetent kids suggests the offense is the primary driver, with empowerment being a defensive reaction rather than a dominant feeling. I should have assigned a significantly higher score to **Offended** (reflecting the personal attack on his work and the implication of his hiding) and a lower score to **Empowered** (reflecting that his "moving forward" is likely a coping mechanism for the offense, not a genuine sense of power). I will adjust the scores to clearly prioritize Offended over Empowered, perhaps setting Empowered lower (e.g., 4) and Offended higher (e.g., 9), while keeping Guilty low as he is not admitting fault.} \\
            }
        }
    }
    & \\ \\

\end{supertabular}
}

\end{document}
